\begin{ex}
Dato l'insieme $A=\{0,1\}$, si considerino i seguenti sottoinsiemi di $A^*$:
\[
X_1=\{01^n001^m\mid n\geq 0,\ m>4\}
\]
\[
X_2=\{a_1\dots a_n\mid n=3k \text{ inizia con } 010 \text{ e termina con } 1011\}
\]
\ \begin{enumerate}
    \item per ognuno dei due insiemi $X_1$ e $X_2$ fornire un automa (deterministico o non deterministico) che lo decide;
    \item costruire a partire dagli automi trovati al punto precedente, l'automa che decide $Y_1=X_1X_2$ e l'automa che decide $Y_2=X_1+X_2$;
    \item fornire l'espressione regolare associata a $Y_1$;
    \item dare la matrice $M$ associata all'automa $Y_2$;
    \item calcolare i primi $3$ termini dell'esecuzione $E(M)$.
\end{enumerate}
\end{ex}

\begin{sol}
\ \begin{enumerate}
    \item L'insieme $X_1$ è deciso dall'automa $\mathcal{A}$:
    \[
    \begin{tikzpicture}[
        scale=0.8, transform shape,
        >=stealth', shorten >=1pt, auto,
        every state/.style={circle, draw, fill=white, minimum size=0.8cm, inner sep=1pt}
    ]
    % Prima riga
    \node[state] (q0) at (0,0)  {$q_0$};
    \node[state] (q1) at (2,0)  {$q_1$};
    \node[state] (q2) at (4,0)  {$q_2$};
    \node[state] (q3) at (6,0)  {$q_3$};
    \node[state] (q4) at (8,0)  {$q_4$};
    \node[state] (q5) at (10,0) {$q_5$};
    % Seconda riga
    \node[state]            (q6) at (10,-2.5) {$q_6$};
    \node[state]            (q7) at (8,-2.5)  {$q_7$};
    \node[state, accepting] (q8) at (6,-2.5)  {$q_8$};

    % Freccia d'ingresso
    \draw[->] ([shift={(-0.7,0.7)}]q0.center) -- (q0);

    % Transizioni
    \draw[->] (q0) -- node {$0$} (q1);
    \draw[->] (q1) to[loop above, looseness=8] node {$1$} (q1);
    \draw[->] (q1) -- node {$0$} (q2);
    \draw[->] (q2) -- node {$0$} (q3);
    \draw[->] (q3) -- node {$1$} (q4);
    \draw[->] (q4) -- node {$1$} (q5);
    \draw[->] (q5) -- node {$1$} (q6);
    \draw[->] (q6) -- node {$1$} (q7);
    \draw[->] (q7) -- node {$1$} (q8);
    \draw[->] (q8) to[loop below, looseness=8] node {$1$} (q8);
    \end{tikzpicture}
    \]
    L'insieme $X_2$ è deciso dall'automa $\mathcal{B}$:
    \[
    \begin{tikzpicture}[
        scale=0.8, transform shape,
        >=stealth', shorten >=1pt, auto,
        every state/.style={circle, draw, fill=white, minimum size=0.8cm, inner sep=1pt}
    ]
    % Prima riga
    \node[state] (q0) at (0,0)    {$q_0$};
    \node[state] (q1) at (2.5,0)  {$q_1$};
    \node[state] (q2) at (5,0)    {$q_2$};
    \node[state] (q3) at (7.5,0)  {$q_3$};
    \node[state] (q4) at (10,0)   {$q_4$};
    % Seconda riga
    \node[state]            (q5) at (10,-2.5)  {$q_5$};
    \node[state]            (q6) at (7.5,-2.5) {$q_6$};
    \node[state]            (q7) at (5,-2.5)   {$q_7$};
    \node[state]            (q8) at (2.5,-2.5) {$q_8$};
    \node[state, accepting] (qF) at (0,-2.5)   {$q_F$};

    % Freccia d'ingresso
    \draw[->] ([shift={(-0.7,0.7)}]q0.center) -- (q0);

    % Transizioni
    \draw[->] (q0) -- node {$0$} (q1);
    \draw[->] (q1) -- node {$1$} (q2);
    \draw[->] (q2) -- node {$0$} (q3);
    \draw[->] (q3) -- node {$0,1$} (q4);
    \draw[->] (q4) -- node {$0,1$} (q5);
    \draw[->] (q5) -- node {$0,1$} (q3);
    \draw[->] (q5) -- node {$1$} (q6);
    \draw[->] (q6) -- node {$0$} (q7);
    \draw[->] (q7) -- node {$1$} (q8);
    \draw[->] (q8) -- node {$1$} (qF);
    \end{tikzpicture}
    \]
    \item L'automa che decide $Y_1=X_1X_2$ si ottiene collegando con una $\varepsilon$-mossa lo stato finale di $\mathcal{A}$ allo stato iniziale di $\mathcal{B}$:
    \[
    \begin{tikzpicture}[
        >=stealth', shorten >=1pt, auto,
        every state/.style={circle, draw, fill=white, minimum size=0.8cm, inner sep=1pt}
    ]
    % Rettangoli di A e B
    \draw[-, line width=2.5pt, rounded corners=3pt] (0,-1.2) rectangle (3.5,1.2);
    \node at (1.75,0) {$\mathcal{A}$};
    \draw[-, line width=2.5pt, rounded corners=3pt] (6.5,-1.2) rectangle (10,1.2);
    \node at (8.25,0) {$\mathcal{B}$};

    % Stati
    \node[state]            (a0) at (0,0)   {$q_0$};
    \node[state]            (aF) at (3.5,0) {$q_F$};
    \node[state]            (b0) at (6.5,0) {$q_0$};
    \node[state, accepting] (bF) at (10,0)  {$q_F$};

    % Freccia d'ingresso e epsilon-mossa
    \draw[->] ([shift={(-0.7,0.7)}]a0.center) -- (a0);
    \draw[->] (aF) -- node {$\varepsilon$} (b0);
    \end{tikzpicture}
    \]
    L'automa che decide $Y_2=X_1+X_2$ si ottiene aggiungendo un nuovo stato iniziale con due $\varepsilon$-mosse verso gli stati iniziali di $\mathcal{A}$ e $\mathcal{B}$:
    \[
    \begin{tikzpicture}[
        >=stealth', shorten >=1pt, auto,
        every state/.style={circle, draw, fill=white, minimum size=0.8cm, inner sep=1pt}
    ]
    % Rettangoli di A e B
    \draw[-, line width=2.5pt, rounded corners=3pt] (3.5,0.5) rectangle (7.5,3.1);
    \node at (5.5,1.8) {$\mathcal{A}$};
    \draw[-, line width=2.5pt, rounded corners=3pt] (3.5,-3.1) rectangle (7.5,-0.5);
    \node at (5.5,-1.8) {$\mathcal{B}$};

    % Stati
    \node[state]            (n0) at (0,0)      {$q_0'$};
    \node[state]            (a0) at (3.5,1.8)  {$q_0$};
    \node[state, accepting] (aF) at (7.5,1.8)  {$q_F$};
    \node[state]            (b0) at (3.5,-1.8) {$q_0$};
    \node[state, accepting] (bF) at (7.5,-1.8) {$q_F$};

    % Freccia d'ingresso e epsilon-mosse
    \draw[->] ([shift={(-0.7,0.7)}]n0.center) -- (n0);
    \draw[->] (n0) -- node {$\varepsilon$} (a0);
    \draw[->] (n0) -- node[swap] {$\varepsilon$} (b0);
    \end{tikzpicture}
    \]
    \item L'espressione regolare associata a $Y_1$ è
    \[
    01^*00\,11111\,1^*\,010\big((0+1)^3\big)^*(0+1)(0+1)1011.
    \]
\end{enumerate}
\end{sol}
\newpage

\begin{ex}
Sia $A=\{0,1\}$.
\ \begin{enumerate}
    \item Dimostrare che l'insieme
    \[
    X=\{ww\mid w\in A^*\}
    \]
    non è DFA-decidibile.
    \item Dimostrare che $X$ è TM-decidibile e che $X\notin \mathrm{DTIME}_1(cn)$.
\end{enumerate}
\end{ex}

\begin{sol}
\ \begin{enumerate}
    \item Supponiamo per assurdo che $X$ sia DFA-decidibile e sia $p>0$ la costante del pumping lemma. Consideriamo
    \[
    w=0^p1^p0^p1^p=(0^p1^p)(0^p1^p)\in X.
    \]
    Sia $w=uxv$ con $|x|>0$ e $|ux|\leq p$. Allora $x=0^t$ con $0<t\leq p$, e quindi per $i=0$
    \[
    ux^0v=0^{p-t}1^p0^p1^p\notin X.
    \]
    Infatti, se fosse $zz$, ognuna delle due metà conterrebbe $p$ uni, quindi $z=0^{p-t}1^p0^s$. Poiché la parola termina con $1$, deve essere $s=0$, da cui $0^p1^p=0^{p-t}1^p$, cioè $t=0$: assurdo.
    \item Usiamo una macchina di Turing a due nastri (bi-infiniti), con $a,b\in\{0,1\}$ nelle etichette. In $q_0$ si copia $w$ sul secondo nastro; in $q_1,q_2$ la testina del secondo nastro si muove il doppio più velocemente di quella del primo, così da portare la prima a metà di $w$ e la seconda all'inizio della copia (se $|w|$ è dispari si rifiuta); in $q_3$ si confrontano le due metà.
    \[
    \begin{tikzpicture}[
        scale=0.8, transform shape,
        >=stealth', shorten >=1pt, auto,
        every state/.style={circle, draw, fill=white, minimum size=0.8cm, inner sep=1pt}
    ]
    % Stati
    \node[state]            (q0)  at (0,0)     {$q_0$};
    \node[state]            (q1)  at (5,0)     {$q_1$};
    \node[state]            (q2)  at (10,0)    {$q_2$};
    \node[state]            (q3)  at (5,-3.5)  {$q_3$};
    \node[state, accepting] (qF1) at (0,-3.5)  {$q_F^1$};
    \node[state, accepting] (qF0) at (10,-3.5) {$q_F^0$};

    % Freccia d'ingresso
    \draw[->] ([shift={(-0.7,0.7)}]q0.center) -- (q0);

    % Copia di w sul secondo nastro
    \draw[->] (q0) to[loop above, looseness=8] node {$0,\square/0,0\;(+1,+1)$} (q0);
    \draw[->] (q0) to[loop below, looseness=8] node {$1,\square/1,1\;(+1,+1)$} (q0);
    \draw[->] (q0) -- node {$\square,\square/\square,\square\;(-1,-1)$} (q1);

    % Ricerca della metà
    \draw[->] (q1) to[bend left=15] node {$a,b/a,b\;(0,-1)$} (q2);
    \draw[->] (q2) to[bend left=15] node {$a,b/a,b\;(-1,-1)$} (q1);
    \draw[->] (q2) -- node {$a,\square/$} (qF0);
    \draw[->] (q1) -- node {$a,\square/a,\square\;(+1,+1)$} (q3);
    \draw[->] (q1) -- node {$\square,\square/$} (qF1);

    % Confronto delle due metà
    \draw[->] (q3) to[loop below, looseness=8] node {$a,a/a,a\;(+1,+1)$} (q3);
    \draw[->] (q3) -- node {$\square,b/$} (qF1);
    \draw[->] (q3) -- node {$a,b/\;(a\neq b)$} (qF0);
    \end{tikzpicture}
    \]
\end{enumerate}
\end{sol}
\newpage

\begin{ex}
Definire a piacere una classe astratta $A$ da cui derivarne una $B$ allo scopo di illustrare i concetti di \emph{inheritance}, \emph{overloading} e di classi astratte.
\end{ex}

\newpage

\begin{ex}
Dato l'automa a stati finiti deterministico $\mathcal{A}$:
\[
\begin{tikzpicture}[
    >=stealth', shorten >=1pt, auto,
    every state/.style={circle, draw, fill=white, minimum size=0.8cm, inner sep=1pt}
]
% Stati
\node[state]            (q0) at (0,0)     {$q_0$};
\node[state]            (q1) at (2.2,0)   {$q_1$};
\node[state, accepting] (q2) at (4.4,0)   {$q_2$};
\node[state]            (q3) at (6.6,0)   {$q_3$};
\node[state, accepting] (q4) at (8.8,0)   {$q_4$};
\node[state]            (q5) at (11,0)    {$q_5$};
\node[state]            (q7) at (7.7,-2.6) {$q_7$};

% Freccia d'ingresso
\draw[->] ([shift={(-0.7,0.7)}]q0.center) -- (q0);

% Transizioni
\draw[->] (q0) to[bend left=20] node {$0$} (q1);
\draw[->] (q1) to[bend left=20] node {$1$} (q0);
\draw[->] (q0) to[loop below, looseness=8] node {$1$} (q0);
\draw[->] (q1) -- node {$0$} (q2);
\draw[->] (q2) to[loop below, looseness=8] node {$1$} (q2);
\draw[->] (q2) -- node {$0$} (q3);
\draw[->] (q3) -- node {$1$} (q4);
\draw[->] (q4) -- node {$0$} (q5);
\draw[->] (q3) -- node[swap] {$0$} (q7);
\draw[->] (q4) -- node {$1$} (q7);
\draw[->] (q5) to[bend right=55] node[swap] {$1$} (q1);
\draw[->] (q5) to[bend right=35] node[swap] {$0$} (q2);
\path[->] (q7) edge[out=200,in=250,looseness=8] node[below left]  {$0$} (q7)
              edge[out=290,in=340,looseness=8] node[below right] {$1$} (q7);
\end{tikzpicture}
\]
\ \begin{enumerate}
    \item calcolare l'espressione regolare associata all'automa $\mathcal{A}$;
    \item a partire dall'espressione regolare trovata al punto precedente, fornire un automa non-deterministico equivalente ad $\mathcal{A}$.
\end{enumerate}
\end{ex}

\begin{ex}
Dato l'automa a stati finiti deterministico $\mathcal{A}$:
\[
\begin{tikzpicture}[
    >=stealth', shorten >=1pt, auto,
    every state/.style={circle, draw, fill=white, minimum size=0.8cm, inner sep=1pt}
]
% Stati
\node[state]            (q0) at (0,0)     {$q_0$};
\node[state]            (q1) at (2.2,0)   {$q_1$};
\node[state, accepting] (q2) at (4.4,0)   {$q_2$};
\node[state]            (q3) at (6.6,0)   {$q_3$};
\node[state, accepting] (q4) at (8.8,0)   {$q_4$};
\node[state]            (q5) at (11,0)    {$q_5$};
\node[state]            (q7) at (7.7,-2.6) {$q_7$};

% Freccia d'ingresso
\draw[->] ([shift={(-0.7,0.7)}]q0.center) -- (q0);

% Transizioni
\draw[->] (q0) to[bend left=20] node {$0$} (q1);
\draw[->] (q1) to[bend left=20] node {$1$} (q0);
\draw[->] (q0) to[loop below, looseness=8] node {$1$} (q0);
\draw[->] (q1) -- node {$0$} (q2);
\draw[->] (q2) to[loop below, looseness=8] node {$1$} (q2);
\draw[->] (q2) -- node {$0$} (q3);
\draw[->] (q3) -- node {$1$} (q4);
\draw[->] (q4) -- node {$0$} (q5);
\draw[->] (q3) -- node[swap] {$0$} (q7);
\draw[->] (q4) -- node {$1$} (q7);
\draw[->] (q5) to[bend right=55] node[swap] {$1$} (q1);
\draw[->] (q5) to[bend right=35] node[swap] {$0$} (q2);
\path[->] (q7) edge[out=200,in=250,looseness=8] node[below left]  {$0$} (q7)
              edge[out=290,in=340,looseness=8] node[below right] {$1$} (q7);
\end{tikzpicture}
\]
\ \begin{enumerate}
    \item calcolare l'espressione regolare associata all'automa $\mathcal{A}$;
    \item a partire dall'espressione regolare trovata al punto precedente, fornire un automa non-deterministico equivalente ad $\mathcal{A}$.
\end{enumerate}
\end{ex}

\begin{sol}
    \
    \begin{enumerate}
    \item Per calcolare l'espressione regolare associata ad $\mathcal{A}$ applichiamo l'algoritmo di eliminazione degli stati, eliminandone uno alla volta. In ogni passo le transizioni nuove sono in grassetto.

    \medskip
    \textbf{Passo 1.} Costruiamo l'automa generalizzato: aggiungiamo un nuovo stato iniziale $q_0^*$ con una $\varepsilon$-mossa verso $q_0$ e un nuovo stato finale $q_F^*$, raggiunto con $\varepsilon$-mosse dagli stati che erano finali ($q_2$ e $q_4$), che smettono di essere finali.
    \[
    \begin{tikzpicture}[
        scale=0.8, transform shape,
        >=stealth', shorten >=1pt, auto,
        every state/.style={circle, draw, fill=white, minimum size=0.8cm, inner sep=1pt}
    ]
    \node[state]            (q0s) at (-2,0)     {$q_0^*$};
    \node[state]            (q0)  at (0,0)      {$q_0$};
    \node[state]            (q1)  at (2.2,0)    {$q_1$};
    \node[state]            (q2)  at (4.4,0)    {$q_2$};
    \node[state]            (q3)  at (6.6,0)    {$q_3$};
    \node[state]            (q4)  at (8.8,0)    {$q_4$};
    \node[state]            (q5)  at (11,0)     {$q_5$};
    \node[state]            (q7)  at (7.7,-2.6) {$q_7$};
    \node[state, accepting] (qF)  at (11,2.5)   {$q_F^*$};

    \draw[->] ([shift={(-0.7,0.7)}]q0s.center) -- (q0s);
    \draw[->, very thick] (q0s) -- node {$\varepsilon$} (q0);
    \draw[->, very thick] (q2) to[bend left=15] node {$\varepsilon$} (qF);
    \draw[->, very thick] (q4) -- node {$\varepsilon$} (qF);

    \draw[->] (q0) to[bend left=20] node {$0$} (q1);
    \draw[->] (q1) to[bend left=20] node {$1$} (q0);
    \draw[->] (q0) to[loop below, looseness=8] node {$1$} (q0);
    \draw[->] (q1) -- node {$0$} (q2);
    \draw[->] (q2) to[loop below, looseness=8] node {$1$} (q2);
    \draw[->] (q2) -- node {$0$} (q3);
    \draw[->] (q3) -- node {$1$} (q4);
    \draw[->] (q4) -- node {$0$} (q5);
    \draw[->] (q3) -- node[swap] {$0$} (q7);
    \draw[->] (q4) -- node {$1$} (q7);
    \draw[->] (q5) to[bend right=55] node[swap] {$1$} (q1);
    \draw[->] (q5) to[bend right=35] node[swap] {$0$} (q2);
    \path[->] (q7) edge[out=200,in=250,looseness=8] node[below left]  {$0$} (q7)
                  edge[out=290,in=340,looseness=8] node[below right] {$1$} (q7);
    \end{tikzpicture}
    \]

    \medskip
    \textbf{Passo 2.} Eliminiamo $q_7$. Da $q_7$ non si raggiunge $q_F^*$ (è uno stato pozzo), quindi nessun cammino accettante passa per $q_7$: si elimina insieme a tutti gli archi che entrano o escono da esso, senza creare transizioni nuove.
    \[
    \begin{tikzpicture}[
        scale=0.8, transform shape,
        >=stealth', shorten >=1pt, auto,
        every state/.style={circle, draw, fill=white, minimum size=0.8cm, inner sep=1pt}
    ]
    \node[state]            (q0s) at (-2,0)   {$q_0^*$};
    \node[state]            (q0)  at (0,0)    {$q_0$};
    \node[state]            (q1)  at (2.2,0)  {$q_1$};
    \node[state]            (q2)  at (4.4,0)  {$q_2$};
    \node[state]            (q3)  at (6.6,0)  {$q_3$};
    \node[state]            (q4)  at (8.8,0)  {$q_4$};
    \node[state]            (q5)  at (11,0)   {$q_5$};
    \node[state, accepting] (qF)  at (11,2.5) {$q_F^*$};

    \draw[->] ([shift={(-0.7,0.7)}]q0s.center) -- (q0s);
    \draw[->] (q0s) -- node {$\varepsilon$} (q0);
    \draw[->] (q2) to[bend left=15] node {$\varepsilon$} (qF);
    \draw[->] (q4) -- node {$\varepsilon$} (qF);

    \draw[->] (q0) to[bend left=20] node {$0$} (q1);
    \draw[->] (q1) to[bend left=20] node {$1$} (q0);
    \draw[->] (q0) to[loop below, looseness=8] node {$1$} (q0);
    \draw[->] (q1) -- node {$0$} (q2);
    \draw[->] (q2) to[loop below, looseness=8] node {$1$} (q2);
    \draw[->] (q2) -- node {$0$} (q3);
    \draw[->] (q3) -- node {$1$} (q4);
    \draw[->] (q4) -- node {$0$} (q5);
    \draw[->] (q5) to[bend right=55] node[swap] {$1$} (q1);
    \draw[->] (q5) to[bend right=35] node[swap] {$0$} (q2);
    \end{tikzpicture}
    \]

    \medskip
    \textbf{Passo 3.} Eliminiamo $q_0$. Gli archi entranti sono $q_0^*\xrightarrow{\varepsilon}q_0$ e $q_1\xrightarrow{1}q_0$, il cappio su $q_0$ è $1$ e l'unico arco uscente è $q_0\xrightarrow{0}q_1$. I nuovi archi sono
    \begin{itemize}
        \item $q_0^*\to q_1$: $\varepsilon\cdot1^*\cdot0=1^*0$;
        \item $q_1\to q_1$: $1\cdot1^*\cdot0=11^*0$.
    \end{itemize}
    \[
    \begin{tikzpicture}[
        scale=0.8, transform shape,
        >=stealth', shorten >=1pt, auto,
        every state/.style={circle, draw, fill=white, minimum size=0.8cm, inner sep=1pt}
    ]
    \node[state]            (q0s) at (0,0)    {$q_0^*$};
    \node[state]            (q1)  at (2.2,0)  {$q_1$};
    \node[state]            (q2)  at (4.4,0)  {$q_2$};
    \node[state]            (q3)  at (6.6,0)  {$q_3$};
    \node[state]            (q4)  at (8.8,0)  {$q_4$};
    \node[state]            (q5)  at (11,0)   {$q_5$};
    \node[state, accepting] (qF)  at (11,2.5) {$q_F^*$};

    \draw[->] ([shift={(-0.7,0.7)}]q0s.center) -- (q0s);
    \draw[->, very thick] (q0s) -- node {$1^*0$} (q1);
    \path[->, very thick] (q1) edge[loop above, looseness=8] node {$11^*0$} (q1);

    \draw[->] (q2) to[bend left=15] node {$\varepsilon$} (qF);
    \draw[->] (q4) -- node {$\varepsilon$} (qF);
    \draw[->] (q1) -- node {$0$} (q2);
    \draw[->] (q2) to[loop below, looseness=8] node {$1$} (q2);
    \draw[->] (q2) -- node {$0$} (q3);
    \draw[->] (q3) -- node {$1$} (q4);
    \draw[->] (q4) -- node {$0$} (q5);
    \draw[->] (q5) to[bend right=55] node[swap] {$1$} (q1);
    \draw[->] (q5) to[bend right=35] node[swap] {$0$} (q2);
    \end{tikzpicture}
    \]

    \medskip
    \textbf{Passo 4.} Eliminiamo $q_3$. L'unico arco entrante è $q_2\xrightarrow{0}q_3$ e l'unico uscente è $q_3\xrightarrow{1}q_4$ (non ci sono cappi). Il nuovo arco è $q_2\to q_4$: $0\cdot1=01$.
    \[
    \begin{tikzpicture}[
        scale=0.8, transform shape,
        >=stealth', shorten >=1pt, auto,
        every state/.style={circle, draw, fill=white, minimum size=0.8cm, inner sep=1pt}
    ]
    \node[state]            (q0s) at (0,0)    {$q_0^*$};
    \node[state]            (q1)  at (2.2,0)  {$q_1$};
    \node[state]            (q2)  at (4.4,0)  {$q_2$};
    \node[state]            (q4)  at (8.8,0)  {$q_4$};
    \node[state]            (q5)  at (11,0)   {$q_5$};
    \node[state, accepting] (qF)  at (11,2.5) {$q_F^*$};

    \draw[->] ([shift={(-0.7,0.7)}]q0s.center) -- (q0s);
    \draw[->] (q0s) -- node {$1^*0$} (q1);
    \path[->] (q1) edge[loop above, looseness=8] node {$11^*0$} (q1);
    \draw[->] (q1) -- node {$0$} (q2);
    \draw[->] (q2) to[loop below, looseness=8] node {$1$} (q2);
    \draw[->, very thick] (q2) -- node {$01$} (q4);
    \draw[->] (q4) -- node {$0$} (q5);
    \draw[->] (q2) to[bend left=15] node {$\varepsilon$} (qF);
    \draw[->] (q4) -- node {$\varepsilon$} (qF);
    \draw[->] (q5) to[bend right=55] node[swap] {$1$} (q1);
    \draw[->] (q5) to[bend right=35] node[swap] {$0$} (q2);
    \end{tikzpicture}
    \]

    \medskip
    \textbf{Passo 5.} Eliminiamo $q_5$. L'unico arco entrante è $q_4\xrightarrow{0}q_5$, gli uscenti sono $q_5\xrightarrow{1}q_1$ e $q_5\xrightarrow{0}q_2$. I nuovi archi sono
    \begin{itemize}
        \item $q_4\to q_1$: $0\cdot1=01$;
        \item $q_4\to q_2$: $0\cdot0=00$.
    \end{itemize}
    \[
    \begin{tikzpicture}[
        scale=0.8, transform shape,
        >=stealth', shorten >=1pt, auto,
        every state/.style={circle, draw, fill=white, minimum size=0.8cm, inner sep=1pt}
    ]
    \node[state]            (q0s) at (0,0)    {$q_0^*$};
    \node[state]            (q1)  at (2.2,0)  {$q_1$};
    \node[state]            (q2)  at (4.4,0)  {$q_2$};
    \node[state]            (q4)  at (8.8,0)  {$q_4$};
    \node[state, accepting] (qF)  at (11,2.5) {$q_F^*$};

    \draw[->] ([shift={(-0.7,0.7)}]q0s.center) -- (q0s);
    \draw[->] (q0s) -- node {$1^*0$} (q1);
    \path[->] (q1) edge[loop above, looseness=8] node {$11^*0$} (q1);
    \draw[->] (q1) -- node {$0$} (q2);
    \draw[->] (q2) to[loop below, looseness=8] node {$1$} (q2);
    \draw[->] (q2) -- node {$01$} (q4);
    \draw[->] (q2) to[bend left=15] node {$\varepsilon$} (qF);
    \draw[->] (q4) -- node {$\varepsilon$} (qF);
    \draw[->, very thick] (q4) to[bend right=55] node[swap] {$01$} (q1);
    \draw[->, very thick] (q4) to[bend right=35] node[swap] {$00$} (q2);
    \end{tikzpicture}
    \]

    \medskip
    \textbf{Passo 6.} Eliminiamo $q_4$. L'unico arco entrante è $q_2\xrightarrow{01}q_4$ e gli uscenti sono $q_4\xrightarrow{\varepsilon}q_F^*$, $q_4\xrightarrow{01}q_1$ e $q_4\xrightarrow{00}q_2$ (non ci sono cappi). I nuovi archi sono
    \begin{itemize}
        \item $q_2\to q_F^*$: $01\cdot\varepsilon=01$;
        \item $q_2\to q_1$: $01\cdot01=0101$;
        \item $q_2\to q_2$ (cappio): $01\cdot00=0100$.
    \end{itemize}
    \[
    \begin{tikzpicture}[
        scale=0.8, transform shape,
        >=stealth', shorten >=1pt, auto,
        every state/.style={circle, draw, fill=white, minimum size=0.8cm, inner sep=1pt}
    ]
    \node[state]            (q0s) at (0,0)    {$q_0^*$};
    \node[state]            (q1)  at (4,0)    {$q_1$};
    \node[state]            (q2)  at (8,0)    {$q_2$};
    \node[state, accepting] (qF)  at (12,0)   {$q_F^*$};

    \draw[->] ([shift={(-0.7,0.7)}]q0s.center) -- (q0s);
    \draw[->] (q0s) -- node {$1^*0$} (q1);
    \path[->] (q1) edge[loop above, looseness=8] node {$11^*0$} (q1);
    \draw[->] (q1) -- node {$0$} (q2);
    \draw[->] (q2) to[loop below, looseness=8] node {$1$} (q2);
    \draw[->] (q2) to[bend left=30] node {$\varepsilon$} (qF);

    \draw[->, very thick] (q2) -- node {$01$} (qF);
    \path[->, very thick] (q2) edge[loop above, looseness=8] node {$0100$} (q2);
    \draw[->, very thick] (q2) to[bend left=30] node {$0101$} (q1);
    \end{tikzpicture}
    \]

    \medskip
    \textbf{Passo 7.} Uniamo gli archi paralleli: i due cappi su $q_2$ diventano un solo cappio $1+0100$ e i due archi da $q_2$ a $q_F^*$ diventano un solo arco $\varepsilon+01$.
    \[
    \begin{tikzpicture}[
        scale=0.8, transform shape,
        >=stealth', shorten >=1pt, auto,
        every state/.style={circle, draw, fill=white, minimum size=0.8cm, inner sep=1pt}
    ]
    \node[state]            (q0s) at (0,0)    {$q_0^*$};
    \node[state]            (q1)  at (4,0)    {$q_1$};
    \node[state]            (q2)  at (8,0)    {$q_2$};
    \node[state, accepting] (qF)  at (12,0)   {$q_F^*$};

    \draw[->] ([shift={(-0.7,0.7)}]q0s.center) -- (q0s);
    \draw[->] (q0s) -- node {$1^*0$} (q1);
    \path[->] (q1) edge[loop above, looseness=8] node {$11^*0$} (q1);
    \draw[->] (q1) -- node {$0$} (q2);
    \draw[->] (q2) to[bend left=30] node {$0101$} (q1);

    \path[->, very thick] (q2) edge[loop below, looseness=8] node {$1+0100$} (q2);
    \draw[->, very thick] (q2) -- node {$\varepsilon+01$} (qF);
    \end{tikzpicture}
    \]

    \medskip
    \textbf{Passo 8.} Eliminiamo $q_2$. L'unico arco entrante è $q_1\xrightarrow{0}q_2$, il cappio è $1+0100$ e gli uscenti sono $q_2\xrightarrow{0101}q_1$ e $q_2\xrightarrow{\varepsilon+01}q_F^*$. I nuovi archi sono
    \begin{itemize}
        \item $q_1\to q_1$ (cappio): $0(1+0100)^*0101$;
        \item $q_1\to q_F^*$: $0(1+0100)^*(\varepsilon+01)$.
    \end{itemize}
    \[
    \begin{tikzpicture}[
        scale=0.8, transform shape,
        >=stealth', shorten >=1pt, auto,
        every state/.style={circle, draw, fill=white, minimum size=0.8cm, inner sep=1pt}
    ]
    \node[state]            (q0s) at (0,0)    {$q_0^*$};
    \node[state]            (q1)  at (5,0)    {$q_1$};
    \node[state, accepting] (qF)  at (11.5,0) {$q_F^*$};

    \draw[->] ([shift={(-0.7,0.7)}]q0s.center) -- (q0s);
    \draw[->] (q0s) -- node {$1^*0$} (q1);
    \draw[->] (q1) to[loop below, looseness=8] node {$11^*0$} (q1);

    \draw[->, very thick] (q1) to[loop above, looseness=8] node {$0(1+0100)^*0101$} (q1);
    \draw[->, very thick] (q1) -- node {$0(1+0100)^*(\varepsilon+01)$} (qF);
    \end{tikzpicture}
    \]

    \medskip
    \textbf{Passo 9.} Eliminiamo $q_1$: i due cappi si uniscono in $11^*0+0(1+0100)^*0101$ e si ottiene la concatenazione dell'arco entrante, del cappio iterato e dell'arco uscente:
    \[
    1^*0\big(11^*0+0(1+0100)^*0101\big)^*0(1+0100)^*(\varepsilon+01).
    \]

    \item Costruiamo un NDFA \emph{senza} $\varepsilon$-mosse seguendo i blocchi dell'espressione
   \[
\begin{tikzpicture}[
    scale=0.85, transform shape,
    >=stealth', shorten >=1pt, auto,
    every state/.style={circle, draw, fill=white, minimum size=0.8cm, inner sep=1pt}
]
% Stati della riga principale
\node[state]            (q0) at (0,0)    {$q_0$};
\node[state]            (q1) at (3,0)    {$q_1$};
\node[state, accepting] (q2) at (6.5,0)  {$q_2$};
\node[state]            (q3) at (9.5,0)  {$q_3$};
\node[state, accepting] (q4) at (12.5,0) {$q_4$};

% Ciclo (b): 11^*0
\node[state] (p1) at (1.5,-3) {$p_1$};

% Ciclo (c): 0(1+0100)^*0101
\node[state] (p2) at (4.8,-3)   {$p_2$};
\node[state] (p3) at (4.8,-4.8) {$p_3$};
\node[state] (p4) at (6.3,-6.2) {$p_4$};
\node[state] (p5) at (7.8,-4.8) {$p_5$};
\node[state] (p6) at (5.8,-2.1) {$p_6$};
\node[state] (p7) at (4.8,-1.3) {$p_7$};
\node[state] (p8) at (3.8,-0.6) {$p_8$};

% Ciclo (d): 0100 su q_2
\node[state] (r1) at (9.5,-3)   {$r_1$};
\node[state] (r2) at (11.8,-3)  {$r_2$};
\node[state] (r3) at (9.5,-1.5) {$r_3$};

% Freccia d'ingresso
\draw[->] ([shift={(-0.7,0.7)}]q0.center) -- (q0);

% Transizioni riga principale
\path[->] (q0) edge[loop below, looseness=8] node {$1$} (q0);
\draw[->] (q0) -- node {$0$} (q1);
\draw[->] (q1) -- node {$0$} (q2);
\path[->] (q2) edge[loop above, looseness=8] node {$1$} (q2);
\draw[->] (q2) -- node {$0$} (q3);
\draw[->] (q3) -- node {$1$} (q4);

% Transizioni ciclo (b)
\draw[->] (q1) to[bend right=15] node[swap] {$1$} (p1);
\path[->] (p1) edge[loop below, looseness=8] node {$1$} (p1);
\draw[->] (p1) to[bend right=15] node[swap] {$0$} (q1);

% Transizioni ciclo (c)
\draw[->] (q1) -- node[swap] {$0$} (p2);
\path[->] (p2) edge[loop left, looseness=8] node {$1$} (p2);
\draw[->] (p2) -- node[swap] {$0$} (p3);
\draw[->] (p3) -- node[swap] {$1$} (p4);
\draw[->] (p4) -- node[swap] {$0$} (p5);
\draw[->] (p5) -- node[swap] {$0$} (p2);
\draw[->] (p2) -- node[swap] {$0$} (p6);
\draw[->] (p6) -- node[swap] {$1$} (p7);
\draw[->] (p7) -- node[swap] {$0$} (p8);
\draw[->] (p8) -- node[swap] {$1$} (q1);

% Transizioni ciclo (d)
\draw[->] (q2) -- node[swap] {$0$} (r1);
\draw[->] (r1) -- node[swap] {$1$} (r2);
\draw[->] (r2) -- node[swap] {$0$} (r3);
\draw[->] (r3) -- node[swap] {$0$} (q2);
\end{tikzpicture}
\]
\end{enumerate}
\end{sol}